\documentclass[11pt]{article}
\usepackage[margin=1in]{geometry}
\usepackage{amsmath,amssymb,amsthm}
\usepackage[T1]{fontenc}
\usepackage{lmodern}
\usepackage{hyperref}

\newtheorem*{theorem}{Theorem}

\begin{document}

\begin{center}
{\Large Literature-implied partial answer: the Laakso-to-tree question for dual Banach spaces}\\[4pt]
{\normalsize Source: M. I. Ostrovskii, arXiv:1102.5082}\\
{\normalsize Supporting paper: M. I. Ostrovskii, arXiv:1302.5968}
\end{center}

\paragraph{Status.}
\textbf{literature\_implied\_answer (partial subcase).}
The implication below settles the source's question for dual Banach spaces,
but not for arbitrary Banach spaces.

\paragraph{Source question.}
In arXiv:1102.5082, Theorem \texttt{T:DiamTree} proves that if the infinite
diamond $D_\omega$ embeds bilipschitzly into a Banach space $X$, then $X$
contains a bounded $\delta$-tree for some $\delta>0$.  Later, in the proof of
Theorem \texttt{T:LaakTree}, the paper proves that bilipschitz embeddability
of the Laakso space $L_\omega$ into $X$ implies that $X$ contains a bounded
$\delta$-semitree and hence fails the Radon--Nikodym property.  At that point
the paper states that it is not known whether bilipschitz embeddability of
$L_\omega$ into $X$ implies the existence of a bounded $\delta$-tree in $X$.

\paragraph{Supporting theorem.}
The later paper arXiv:1302.5968 proves in Theorem 1.10 that a dual Banach
space does not have the Radon--Nikodym property if and only if it admits a
bilipschitz embedding of $D_\omega$.

\begin{theorem}
Let $X$ be a dual Banach space.  If $L_\omega$ embeds bilipschitzly into $X$,
then $X$ contains a bounded $\delta$-tree for some $\delta>0$.
\end{theorem}

\begin{proof}
Assume that $L_\omega$ embeds bilipschitzly into the dual Banach space $X$.
By arXiv:1102.5082, Theorem \texttt{T:LaakTree}, the space $X$ does not have
the Radon--Nikodym property.  Since $X$ is a dual Banach space, arXiv:1302.5968,
Theorem 1.10, implies that $D_\omega$ embeds bilipschitzly into $X$.  Applying
arXiv:1102.5082, Theorem \texttt{T:DiamTree}, to this embedding of
$D_\omega$, we obtain a bounded $\delta$-tree in $X$ for some $\delta>0$.
\end{proof}

\paragraph{Scope and provenance.}
This is not a full answer to the source question, because the dual-space
hypothesis is essential to the supporting theorem used here.  The relation is
an agent-identified implication from a later theorem; arXiv:1302.5968 does not
explicitly state that it answers the sentence in arXiv:1102.5082.

\begin{thebibliography}{9}

\bibitem{Ost11}
M. I. Ostrovskii,
\emph{On metric characterizations of some classes of Banach spaces},
arXiv:1102.5082.

\bibitem{Ost13}
M. I. Ostrovskii,
\emph{On metric characterizations of the Radon-Nikodym and related properties
of Banach spaces},
arXiv:1302.5968.

\end{thebibliography}

\end{document}
